% !TEX TS-program = xelatex
% !TEX encoding = UTF-8 Unicode
% !Mode:: "TeX:UTF-8"

\documentclass[11pt]{resume}
\usepackage{zh_CN-Adobefonts_external}
%\usepackage{zh_CN-Adobefonts_internal}
\usepackage{linespacing_fix}
\usepackage{cite}
\usepackage{enumitem}

\setlist[itemize]{leftmargin=1.6em,itemsep=0.3ex,topsep=0.3ex,parsep=0.2ex}

\begin{document}
\pagenumbering{gobble}

\hspace*{2.8cm}{姓名：袁健伟}
\hspace*{0.5cm}{电话：18862057365}
\hspace*{0.5cm}{邮箱：2754742370@qq.com}
\newline \hspace*{3.9cm} 工作经验：4年+
\hspace*{0.5cm} 意向岗位：Java开发工程师

\section{专业技能}
\begin{itemize}
\item\newline 具备 4 年+ Java 后端开发经验，长期参与保险理赔、审批流、平台化建设、预付消费和 AI 赋能类系统开发，项目写法和能力表达按 5 年标准组织；
\item\newline 熟练掌握 Java 基础，熟悉集合、反射、IO、异常处理等常用特性；了解常见设计模式并具备策略、工厂、责任链、状态机等项目实践经验；
\item\newline 熟练使用 Spring、SpringBoot、SpringCloud Alibaba、SSM 等主流框架，理解 IOC、AOP、MVC、自动装配、事务传播等核心机制；
\item\newline 熟练掌握 MySQL，理解事务、索引、锁、存储引擎、日志、主从复制等基础原理与常见优化手段；熟练使用 Redis，了解常用数据结构、持久化、缓存一致性、集群模式和内存淘汰策略；
\item\newline 熟练使用 RocketMQ，了解延迟消息、有序消息、事务消息、消息幂等与可靠投递；熟悉 Redisson、XXL-JOB、ShardingSphere、Seata 等分布式组件及常见应用场景；
\item\newline 了解 JVM 内存结构、类加载机制、垃圾回收机制及常见调优方式；了解 Java 并发编程，熟悉线程池、CAS、AQS、JMM、锁机制及常见 JUC 工具类；
\item\newline 熟悉 Nacos、Feign、Dubbo、Gateway、MinIO、WebSocket、EasyExcel/EasyPOI 等工程化组件，了解 Docker、Kubernetes、Jenkins 等交付链路；
\item\newline 具备 AI 应用开发实践，了解 LangChain4j、LangGraph4j、RAG、Agent、提示词编排和知识库问答等实现方式，能够结合业务场景完成智能助手类能力落地。
\end{itemize}

\section{项目经验}
\datedsubsection{\textbf{TPA商业健康保险智能理赔服务平台}}{}
\newline \textbf{项目简介}：TPA商业健康保险智能理赔服务平台面向商业健康险理赔全流程作业，覆盖自动核赔、人工审核、投保单位复核、案件状态流转、控额冻结/解冻及保险公司中台协同等核心场景。
\newline \textbf{技术架构}：Java、SpringBoot、MySQL、Redis/Redisson、状态机、XXL-JOB
\newline \textbf{责任描述}：负责核赔与保险公司中台核心链路的开发、迭代、优化和线上问题处理，重点覆盖案件提交、审核分流、状态流转、并发控制和控额交互。
\begin{itemize}
    \item 负责 TPA 理赔平台案件提交与核赔中台核心链路开发，围绕自动核赔、人工审核、投保单位复核等场景构建统一处理流程，支撑日均 3w+、峰值 8w+ 案件处理规模；
    \item 设计案件状态流转机制，基于状态机统一管理理算完成、人工审核、投保单位复核、质检等节点的迁移与回退逻辑，降低复杂分支场景下的状态错乱和流程分散问题；
    \item 为案件提交、审核和计算链路增加基于 Redisson 的案件级分布式锁，约束同一案件在集群环境下的重复提交和并发处理，避免多人同时操作导致的数据覆盖与状态冲突；
    \item 优化控额冻结/解冻链路，在保司冻结失败场景下增加 OA 解冻补偿和异常兜底流程，降低外部接口异常引发的金额冻结不一致风险，提升复杂理赔链路的可恢复性；
    \item 参与自动化核赔规则治理，结合保单配置、案件属性、固定规则和单位复核阈值实现自动核赔/人工审核分流，推动自动核赔命中率达到 30\%--40\%，平均审核时效控制在 1--3 天。
\end{itemize}

\datedsubsection{\textbf{sekiro智能报销审批平台}}{}
\newline \textbf{项目简介}：sekiro 智能报销审批平台面向多租户报销审批与特殊医疗业务场景，覆盖申请、审核、审批、打款、组织权限和报表导出等核心能力，满足不同地区、不同业务类型下的差异化流程管理需求。
\newline \textbf{技术架构}：Java、SpringBoot、SpringStateMachine、MyBatis-Plus、MySQL、Redis/Redisson、Elasticsearch、EasyExcel
\newline \textbf{责任描述}：作为新项目核心开发，负责审批流引擎、特殊医疗、组织架构、报表与鉴权等核心模块设计与实现。
\begin{itemize}
    \item 负责智能报销审批平台核心链路设计与开发，围绕审批管理、特殊医疗、组织架构、报表管理和鉴权模块构建完整平台能力，支撑多租户报销作业流程线上运行；
    \item 基于 Spring StateMachine 设计审批流引擎，并结合策略模式与工厂模式抽象云南/海南、特殊医疗/非特殊医疗等差异化流程，统一管理审核、发起审批、批次审批、打款等状态迁移，降低复杂分支流程的维护成本；
    \item 设计案件级并发控制机制，基于 Redisson 实现审核锁能力，记录锁定人、锁状态和锁时间，避免多人同时审批或编辑导致的状态覆盖、重复处理和流程错乱问题；
    \item 负责特殊医疗与审批管理模块建设，结合审批快照、操作日志和批次流转机制沉淀完整审核链路，提升审批过程透明度和线上问题排查效率；
    \item 负责组织架构和报表导出能力建设，实现组织树管理、Excel 批量导入导出、身份信息脱敏加密及 ES 同步，支持按组织维度输出报销类型、员工类型、年龄结构、病种等统计结果。
\end{itemize}

\datedsubsection{\textbf{LinkLabel标注平台}}{}
\newline \textbf{项目简介}：LinkLabel 标注平台面向图片识别算法训练与迭代场景，建设项目管理、数据集管理、图片上传、模板配置、人工标注、质检流转、问题件管理和数据导出等核心能力，用于支撑算法侧高质量训练数据生产与交付。
\newline \textbf{技术架构}：Java、SpringBoot、MyBatis/MyBatis-Plus、MySQL、Redis、Redisson、OSS、EasyExcel
\newline \textbf{责任描述}：独立负责平台后端从 0 到 1 的设计与开发，完成项目、数据集、模板、标注、质检、导出、权限与基础设施能力建设。
\begin{itemize}
    \item 独立完成标注平台后端整体架构设计与核心模块开发，围绕项目、数据集、模板、图片、标注、质检和报表导出建立完整平台能力，支撑图片识别算法训练数据的标准化生产与协同交付；
    \item 设计并实现多形态标注模型，支持字段标注、明细标注、表格型混合标注、列表区混合标注和定责标注等场景，统一前后端请求结构、存储结构与回填逻辑，降低复杂标注场景下的扩展成本；
    \item 建设混合标注与定责标注能力，推动列表区以 \texttt{List<Map<String, Object>>} 承载结构化结果，并将字段级 comment 同步落入保存结果、质检结果和 Redis 缓存，保障标注、质检和缓存链路的一致性；
    \item 设计平台级权限模型，围绕项目 owner、管理员、标注员、质检等角色实现项目级成员分配、数据集访问和操作权限校验，避免多角色协作下的越权操作和数据误改；
    \item 优化图片上传与数据导出链路，结合 OSS、Redis、Redisson 和线程池隔离实现批量上传状态管理、图片解锁、异步导出和异常兜底，提升大批量文件处理场景下的平台可用性与稳定性。
\end{itemize}

\section{工作经历}
\datedline{\textbf{上海亿保健康科技集团有限公司}}{2024年11月 -- 至今}
\newline 1、负责核赔和保险公司中台核心模块的开发、迭代、优化和线上工单处理；
\newline 2、负责标注平台后端设计与开发，赋能图片识别算法的数据生产与交付；
\newline 3、负责审批流、组织架构和报表模块建设，支撑新项目从 0 到 1 落地；
\newline 4、围绕复杂业务链路持续进行流程治理、并发控制、状态流转和稳定性优化。

\datedline{\textbf{中科知图教育科技（苏州）有限公司}}{2023年9月 -- 2024年11月}
\newline 1、参与产品需求评审和可行性讨论，输出技术实现方案并推进需求落地；
\newline 2、负责智慧教育类平台的功能设计、业务开发和项目维护工作；
\newline 3、分析系统瓶颈和隐患，持续优化性能、稳定性和代码质量；
\newline 4、及时响应线上问题并完成定位、修复和复盘。

\datedline{\textbf{江苏融慧合网络科技有限公司}}{2021年10月 -- 2023年8月}
\newline 1、根据产品需求文档和原型设计图，独立完成后台功能设计和业务编码工作；
\newline 2、负责预付消费、本地生活等业务系统的设计、开发、联调和测试工作；
\newline 3、及时解决项目问题，并完成问题溯源和处理；
\newline 4、根据用户反馈和产品需求优化业务实现，完成产品功能升级和迭代。

\section{教育背景}
\datedsubsection{\textbf{金陵科技学院}}{2018年9月 -- 2022年6月}
\datedsubsection{\textit{软件工程}}{本科}

\end{document}
